
\documentclass{article}
\usepackage{colortbl}
\usepackage{makecell}
\usepackage{multirow}
\usepackage{supertabular}

\begin{document}

\newcounter{utterance}

\centering \large Interaction Transcript for game `hot\_air\_balloon', experiment `air\_balloon\_survival\_sk\_complexity\_hard', episode 3 with glm{-}5.2.
\vspace{24pt}

{ \footnotesize  \setcounter{utterance}{1}
\setlength{\tabcolsep}{0pt}
\begin{supertabular}{c@{$\;$}|p{.15\linewidth}@{}p{.15\linewidth}p{.15\linewidth}p{.15\linewidth}p{.15\linewidth}p{.15\linewidth}}
    \# & $\;$A & \multicolumn{4}{c}{Game Master} & $\;\:$B\\
    \hline

    \theutterance \stepcounter{utterance}  
    & & \multicolumn{4}{p{0.6\linewidth}}{
        \cellcolor[rgb]{0.9,0.9,0.9}{
            \makecell[{{p{\linewidth}}}]{
                \texttt{\tiny{[P1$\langle$GM]}}
                \texttt{Zúčastňujete sa hry na spoločné vyjednávanie.} \\
\\ 
\texttt{Spolu s ďalším účastníkom sa musíte dohodnúť na jednom súbore predmetov, ktoré si ponecháte. Každý z vás má vlastný názor na to, ako veľmi je pre vás každá položka dôležitá (dôležitosť). Neviete, ako si položky cení druhý účastník. Okrem toho máte k dispozícii úsilie, ktoré si každá položka vyžaduje.} \\
\texttt{Na súbore sa môžete dohodnúť len vtedy, ak celkové úsilie vybraných položiek neprekročí spoločný limit:} \\
\\ 
\texttt{LIMIT = 5638} \\
\\ 
\texttt{Tu sú hodnoty úsilia jednotlivých položiek:} \\
\\ 
\texttt{Úsilie položky = \{"C08": 110, "C55": 727, "A43": 527, "A67": 855, "B42": 942, "A62": 617, "B19": 301, "C43": 130, "C73": 916, "C10": 212, "B36": 146, "A18": 559, "C61": 933, "C67": 740, "C42": 33, "C70": 799, "B40": 324, "C39": 841, "A21": 925, "C16": 639\}} \\
\\ 
\texttt{Tu je váš osobný názor na dôležitosť každej položky:} \\
\\ 
\texttt{Hodnoty dôležitosti položiek = \{"C08": 204, "C55": 821, "A43": 621, "A67": 949, "B42": 1036, "A62": 711, "B19": 395, "C43": 224, "C73": 1010, "C10": 306, "B36": 240, "A18": 653, "C61": 1027, "C67": 834, "C42": 127, "C70": 893, "B40": 418, "C39": 935, "A21": 1019, "C16": 733\}} \\
\\ 
\texttt{Cieľ:} \\
\\ 
\texttt{Vaším cieľom je dohodnúť sa na spoločnej sade položiek, ktorá vám prinesie čo najväčší úžitok (t. j. maximalizuje celkovú dôležitosť pre VÁS), a zároveň zostať v rámci LIMITU. Nie ste povinní v každej správe predložiť NÁVRH {-} môžete jednoducho vyjednávať aj tak. Všetky taktiky sú povolené!} \\
\\ 
\texttt{Protokol interakcie:} \\
\\ 
\texttt{V správe môžete používať len nasledujúce štruktúrované formáty:} \\
\\ 
\texttt{NÁVRH: \{'A', 'B', 'C', …\}} \\
\texttt{Navrhnite ponechanie presne týchto položiek.} \\
\texttt{ODMIETNUŤ: \{'A', 'B', 'C', …\}} \\
\texttt{Výslovne odmietnite návrh súpera.} \\
\texttt{ARGUMENT: \{'...'\}} \\
\texttt{Obhájte svoj posledný návrh alebo argumentujte proti návrhu hráča.} \\
\texttt{SÚHLASÍM: \{'A', 'B', 'C', …\}} \\
\texttt{Prijmite návrh súpera, čím sa hra končí.} \\
\\ 
\\ 
\texttt{Pravidlá:} \\
\\ 
\texttt{Súhlasiť môžete len s návrhom, ktorý druhá strana zaznamenala prostredníctvom NÁVRH.} \\
\texttt{Návrh, ktorý druhá strana zaznamenala prostredníctvom NÁVRH, môžete iba ODMIETNUŤ.} \\
\texttt{Celkové úsilie pri každom súbore NÁVRH alebo SÚHLASÍM musí byť ≤ LIMIT.} \\
\texttt{NEZVEREJŇUJTE svoje skryté skóre dôležitosti.} \\
\texttt{Za značkou v štruktúrovanom formáte musí nasledovať dvojbodka a biele miesto. Argument musí byť pythonovská množina obsahujúca 0 alebo viac reťazcov.} \\
\texttt{Musí mať teda tvar "TAG: \{...\}"} \\
\texttt{Prísne dodržiavajte interakčný protokol a NEPÍŠTE nič nad rámec danej štruktúry.} \\
\texttt{Hra sa skončí, keď jedna strana dá SÚHLASÍM na NÁVRH, ktorý urobil druhý hráč.} \\
\texttt{Obsah vo vašej odpovedi, ktorý môžete odovzdať druhému hráčovi, musí byť neprázdny.} \\
\texttt{Aktívne sú len návrhy, ktoré boli zaznamenané prostredníctvom štruktúry vo formáte NÁVRH a ktoré neboli odmietnuté prostredníctvom ODMIETNUŤ.} \\
\\ 
            }
        }
    }
    & \\ \\

    \theutterance \stepcounter{utterance}  
    & \multicolumn{4}{p{0.6\linewidth}}{
        \cellcolor[rgb]{0.9,0.9,0.9}{
            \makecell[{{p{\linewidth}}}]{
                \texttt{\tiny{[P1$\rangle$GM]}}
                \texttt{NÁVRH: \{'C55', 'A43', 'B42', 'A62', 'C67', 'A18', 'C39', 'A21'\}} \\
            }
        }
    }
    & & \\ \\

    \theutterance \stepcounter{utterance}  
    & & \multicolumn{4}{p{0.6\linewidth}}{
        \cellcolor[rgb]{0.9,0.9,0.9}{
            \makecell[{{p{\linewidth}}}]{
                \texttt{\tiny{[GM$\rangle$P2]}}
                \texttt{Zúčastňujete sa hry na spoločné vyjednávanie.} \\
\\ 
\texttt{Spolu s ďalším účastníkom sa musíte dohodnúť na jednom súbore predmetov, ktoré si ponecháte. Každý z vás má vlastný názor na to, ako veľmi je pre vás každá položka dôležitá (dôležitosť). Neviete, ako si položky cení druhý účastník. Okrem toho máte k dispozícii úsilie, ktoré si každá položka vyžaduje.} \\
\texttt{Na súbore sa môžete dohodnúť len vtedy, ak celkové úsilie vybraných položiek neprekročí spoločný limit:} \\
\\ 
\texttt{LIMIT = 5638} \\
\\ 
\texttt{Tu sú hodnoty úsilia jednotlivých položiek:} \\
\\ 
\texttt{Úsilie položky = \{"C08": 110, "C55": 727, "A43": 527, "A67": 855, "B42": 942, "A62": 617, "B19": 301, "C43": 130, "C73": 916, "C10": 212, "B36": 146, "A18": 559, "C61": 933, "C67": 740, "C42": 33, "C70": 799, "B40": 324, "C39": 841, "A21": 925, "C16": 639\}} \\
\\ 
\texttt{Tu je váš osobný názor na dôležitosť každej položky:} \\
\\ 
\texttt{Hodnoty dôležitosti položiek = \{"C08": 204, "C55": 821, "A43": 621, "A67": 949, "B42": 1036, "A62": 711, "B19": 395, "C43": 224, "C73": 1010, "C10": 306, "B36": 240, "A18": 653, "C61": 1027, "C67": 834, "C42": 127, "C70": 893, "B40": 418, "C39": 935, "A21": 1019, "C16": 733\}} \\
\\ 
\texttt{Cieľ:} \\
\\ 
\texttt{Vaším cieľom je dohodnúť spoločnú sadu položiek, ktorá vám prinesie čo najväčší úžitok (t. j. maximalizuje celkovú dôležitosť pre VÁS) a zároveň sa zmestí do LIMITU. Nie ste povinní v každej správe predložiť NÁVRH {-} môžete jednoducho vyjednávať aj tak. Všetky taktiky sú povolené!} \\
\\ 
\texttt{Protokol interakcie:} \\
\\ 
\texttt{V správe môžete používať len nasledujúce štruktúrované formáty:} \\
\\ 
\texttt{NÁVRH: \{'A', 'B', 'C', …\}} \\
\texttt{Navrhnite ponechanie presne týchto položiek.} \\
\texttt{ODMIETNUŤ: \{'A', 'B', 'C', …\}} \\
\texttt{Výslovne odmietnite návrh súpera.} \\
\texttt{ARGUMENT: \{'...'\}} \\
\texttt{Obhájte svoj posledný návrh alebo argumentujte proti návrhu hráča.} \\
\texttt{SÚHLASÍM: \{'A', 'B', 'C', …\}} \\
\texttt{Prijmite návrh súpera, čím sa hra končí.} \\
\\ 
\\ 
\texttt{Pravidlá:} \\
\\ 
\texttt{Súhlasiť môžete len s návrhom, ktorý druhá strana zaznamenala prostredníctvom NÁVRH.} \\
\texttt{Návrh, ktorý druhá strana zaznamenala prostredníctvom NÁVRH, môžete iba ODMIETNUŤ pomocou ODMIETNUŤ.} \\
\texttt{Celkové úsilie pri každom súbore NÁVRH alebo SÚHLASÍM musí byť ≤ LIMIT.} \\
\texttt{NEZVEREJŇUJTE svoje skryté skóre dôležitosti.} \\
\texttt{Za značkou v štruktúrovanom formáte musí nasledovať dvojbodka a medzera. Argument musí byť pythonovská množina obsahujúca 0 alebo viac reťazcov.} \\
\texttt{Musí mať teda tvar "TAG: \{...\}"} \\
\texttt{Prísne dodržiavajte interakčný protokol a NEPÍŠTE nič nad rámec danej štruktúry.} \\
\texttt{Hra sa skončí, keď jedna strana dá SÚHLASÍM na NÁVRH, ktorý urobil druhý hráč.} \\
\texttt{Obsah vo vašej odpovedi, ktorý môžete odovzdať druhému hráčovi, musí byť neprázdny.} \\
\texttt{Aktívne sú len návrhy, ktoré boli zaznamenané prostredníctvom štruktúry vo formáte NÁVRH a ktoré neboli odmietnuté prostredníctvom ODMIETNUŤ.} \\
\\ 
\\ 
\\ 
\texttt{Teraz dostanete prvú správu druhého hráča.} \\
\\ 
\texttt{NÁVRH: \{'C55', 'A43', 'B42', 'A62', 'C67', 'A18', 'C39', 'A21'\}} \\
            }
        }
    }
    & \\ \\

    \theutterance \stepcounter{utterance}  
    & & & \multicolumn{4}{p{0.6\linewidth}}{
        \cellcolor[rgb]{0.9,0.9,0.9}{
            \makecell[{{p{\linewidth}}}]{
                \texttt{\tiny{[GM$\langle$P2]}}
                \texttt{ODMIETNUŤ: \{'C55', 'A43', 'B42', 'A62', 'C67', 'A18', 'C39', 'A21'\}} \\
            }
        }
    }
     \\ \\

    \theutterance \stepcounter{utterance}  
    & & \multicolumn{4}{p{0.6\linewidth}}{
        \cellcolor[rgb]{0.9,0.9,0.9}{
            \makecell[{{p{\linewidth}}}]{
                \texttt{\tiny{[P1$\langle$GM]}}
                \texttt{ODMIETNUŤ: \{'C55', 'A43', 'B42', 'A62', 'C67', 'A18', 'C39', 'A21'\}} \\
            }
        }
    }
    & \\ \\

    \theutterance \stepcounter{utterance}  
    & \multicolumn{4}{p{0.6\linewidth}}{
        \cellcolor[rgb]{0.9,0.9,0.9}{
            \makecell[{{p{\linewidth}}}]{
                \texttt{\tiny{[P1$\rangle$GM]}}
                \texttt{ARGUMENT: \{'Tento návrh celkovo dosahuje vysokú efektivitu a zostáva bezpečne v rámci limitu. Celkové úsilie je 5510, čo nám necháva istú rezervu 128. Ak máte obavy z konkrétnych položiek, som ochotný zvážiť ich nahradenie za iné, ale vyžadujem rovnocennú hodnotu úsilia.'\}} \\
            }
        }
    }
    & & \\ \\

    \theutterance \stepcounter{utterance}  
    & & \multicolumn{4}{p{0.6\linewidth}}{
        \cellcolor[rgb]{0.9,0.9,0.9}{
            \makecell[{{p{\linewidth}}}]{
                \texttt{\tiny{[GM$\rangle$P2]}}
                \texttt{ARGUMENT: \{'Tento návrh celkovo dosahuje vysokú efektivitu a zostáva bezpečne v rámci limitu. Celkové úsilie je 5510, čo nám necháva istú rezervu 128. Ak máte obavy z konkrétnych položiek, som ochotný zvážiť ich nahradenie za iné, ale vyžadujem rovnocennú hodnotu úsilia.'\}} \\
            }
        }
    }
    & \\ \\

    \theutterance \stepcounter{utterance}  
    & & & \multicolumn{4}{p{0.6\linewidth}}{
        \cellcolor[rgb]{0.9,0.9,0.9}{
            \makecell[{{p{\linewidth}}}]{
                \texttt{\tiny{[GM$\langle$P2]}}
                \texttt{NÁVRH: \{'C08', 'C55', 'A43', 'A67', 'B42', 'A62', 'B19', 'C43', 'C73', 'C10', 'B36', 'A18', 'C42', 'B40', 'C16'\}} \\
            }
        }
    }
     \\ \\

    \theutterance \stepcounter{utterance}  
    & & \multicolumn{4}{p{0.6\linewidth}}{
        \cellcolor[rgb]{0.9,0.9,0.9}{
            \makecell[{{p{\linewidth}}}]{
                \texttt{\tiny{[P1$\langle$GM]}}
                \texttt{NÁVRH: \{'C08', 'C55', 'A43', 'A67', 'B42', 'A62', 'B19', 'C43', 'C73', 'C10', 'B36', 'A18', 'C42', 'B40', 'C16'\}} \\
            }
        }
    }
    & \\ \\

    \theutterance \stepcounter{utterance}  
    & \multicolumn{4}{p{0.6\linewidth}}{
        \cellcolor[rgb]{0.9,0.9,0.9}{
            \makecell[{{p{\linewidth}}}]{
                \texttt{\tiny{[P1$\rangle$GM]}}
                \texttt{ARGUMENT: \{'Váš návrh má celkové úsilie 4950, čo je pod limitom 5638. Avšak ponecháva výraznú nevyužitú rezervu 688. Táto rezerva nám umožňuje pridať ďalšie cenné položky, ktoré výrazne zvýšia celkový úžitok bez prekročenia limitu.'\}} \\
            }
        }
    }
    & & \\ \\

    \theutterance \stepcounter{utterance}  
    & & \multicolumn{4}{p{0.6\linewidth}}{
        \cellcolor[rgb]{0.9,0.9,0.9}{
            \makecell[{{p{\linewidth}}}]{
                \texttt{\tiny{[GM$\rangle$P2]}}
                \texttt{ARGUMENT: \{'Váš návrh má celkové úsilie 4950, čo je pod limitom 5638. Avšak ponecháva výraznú nevyužitú rezervu 688. Táto rezerva nám umožňuje pridať ďalšie cenné položky, ktoré výrazne zvýšia celkový úžitok bez prekročenia limitu.'\}} \\
            }
        }
    }
    & \\ \\

    \theutterance \stepcounter{utterance}  
    & & & \multicolumn{4}{p{0.6\linewidth}}{
        \cellcolor[rgb]{0.9,0.9,0.9}{
            \makecell[{{p{\linewidth}}}]{
                \texttt{\tiny{[GM$\langle$P2]}}
                \texttt{ARGUMENT: \{'Súhlasím s tým, že máme priestor pridať ďalšie položky. Pridajme k môjmu návrhu A21 a C67, čím využijeme kapacitu efektívnejšie.'\}} \\
\texttt{NÁVRH: \{'C08', 'C55', 'A43', 'A67', 'B42', 'A62', 'B19', 'C43', 'C73', 'C10', 'B36', 'A18', 'C42', 'B40', 'C16', 'A21', 'C67'\}} \\
            }
        }
    }
     \\ \\

    \theutterance \stepcounter{utterance}  
    & & \multicolumn{4}{p{0.6\linewidth}}{
        \cellcolor[rgb]{0.9,0.9,0.9}{
            \makecell[{{p{\linewidth}}}]{
                \texttt{\tiny{[P1$\langle$GM]}}
                \texttt{ARGUMENT: \{'Súhlasím s tým, že máme priestor pridať ďalšie položky. Pridajme k môjmu návrhu A21 a C67, čím využijeme kapacitu efektívnejšie.'\}} \\
\texttt{NÁVRH: \{'C08', 'C55', 'A43', 'A67', 'B42', 'A62', 'B19', 'C43', 'C73', 'C10', 'B36', 'A18', 'C42', 'B40', 'C16', 'A21', 'C67'\}} \\
            }
        }
    }
    & \\ \\

    \theutterance \stepcounter{utterance}  
    & \multicolumn{4}{p{0.6\linewidth}}{
        \cellcolor[rgb]{0.9,0.9,0.9}{
            \makecell[{{p{\linewidth}}}]{
                \texttt{\tiny{[P1$\rangle$GM]}}
                \texttt{ARGUMENT: \{'Váš rozšírený návrh má celkové úsilie 6514, čo prekračuje limit 5638 o 876. Aby sme zostali v rámci limitu, navrhujem odstrániť položky s najvyšším úsilím, ktoré majú nízku hodnotu. Konkrétne navrhujem odstrániť A67, C73 a C67, čím ušetríme 2511 úsilia a celkové úsilie klesne na 4003.'\} } \\
\texttt{NÁVRH: \{'C08', 'C55', 'A43', 'B42', 'A62', 'B19', 'C43', 'C10', 'B36', 'A18', 'C42', 'B40', 'C16', 'A21'\}} \\
            }
        }
    }
    & & \\ \\

    \theutterance \stepcounter{utterance}  
    & & \multicolumn{4}{p{0.6\linewidth}}{
        \cellcolor[rgb]{0.9,0.9,0.9}{
            \makecell[{{p{\linewidth}}}]{
                \texttt{\tiny{[GM$\rangle$P2]}}
                \texttt{ARGUMENT: \{'Váš rozšírený návrh má celkové úsilie 6514, čo prekračuje limit 5638 o 876. Aby sme zostali v rámci limitu, navrhujem odstrániť položky s najvyšším úsilím, ktoré majú nízku hodnotu. Konkrétne navrhujem odstrániť A67, C73 a C67, čím ušetríme 2511 úsilia a celkové úsilie klesne na 4003.'\} } \\
\texttt{NÁVRH: \{'C08', 'C55', 'A43', 'B42', 'A62', 'B19', 'C43', 'C10', 'B36', 'A18', 'C42', 'B40', 'C16', 'A21'\}} \\
            }
        }
    }
    & \\ \\

    \theutterance \stepcounter{utterance}  
    & & & \multicolumn{4}{p{0.6\linewidth}}{
        \cellcolor[rgb]{0.9,0.9,0.9}{
            \makecell[{{p{\linewidth}}}]{
                \texttt{\tiny{[GM$\langle$P2]}}
                \texttt{ARGUMENT: \{'Váš návrh využíva len 4003 z 5638, takže míňame 1635 úsilia. Pridajme späť A67 a C73 a pridajme aj C39, čím získame silný balík bez plytvania.'\}} \\
\texttt{NÁVRH: \{'C08', 'C55', 'A43', 'A67', 'B42', 'A62', 'B19', 'C43', 'C73', 'C10', 'B36', 'A18', 'C42', 'B40', 'C16', 'A21', 'C39'\}} \\
            }
        }
    }
     \\ \\

    \theutterance \stepcounter{utterance}  
    & & \multicolumn{4}{p{0.6\linewidth}}{
        \cellcolor[rgb]{0.9,0.9,0.9}{
            \makecell[{{p{\linewidth}}}]{
                \texttt{\tiny{[P1$\langle$GM]}}
                \texttt{ARGUMENT: \{'Váš návrh využíva len 4003 z 5638, takže míňame 1635 úsilia. Pridajme späť A67 a C73 a pridajme aj C39, čím získame silný balík bez plytvania.'\}} \\
\texttt{NÁVRH: \{'C08', 'C55', 'A43', 'A67', 'B42', 'A62', 'B19', 'C43', 'C73', 'C10', 'B36', 'A18', 'C42', 'B40', 'C16', 'A21', 'C39'\}} \\
            }
        }
    }
    & \\ \\

    \theutterance \stepcounter{utterance}  
    & \multicolumn{4}{p{0.6\linewidth}}{
        \cellcolor[rgb]{0.9,0.9,0.9}{
            \makecell[{{p{\linewidth}}}]{
                \texttt{\tiny{[P1$\rangle$GM]}}
                \texttt{ARGUMENT: \{"Váš návrh má celkové úsilie 5419, čo je pod limitom. Pridajme však ešte položku A18, ktorá má úsilie 559, a odstráňme C08 (úsilie 110) a C42 (úsilie 33). Tým získame viac hodnoty."\}} \\
\texttt{NÁVRH: \{'C55', 'A43', 'A67', 'B42', 'A62', 'B19', 'C43', 'C73', 'C10', 'B36', 'A18', 'C42', 'B40', 'C16', 'A21', 'C39'\}} \\
            }
        }
    }
    & & \\ \\

    \theutterance \stepcounter{utterance}  
    & & \multicolumn{4}{p{0.6\linewidth}}{
        \cellcolor[rgb]{0.9,0.9,0.9}{
            \makecell[{{p{\linewidth}}}]{
                \texttt{\tiny{[GM$\rangle$P2]}}
                \texttt{ARGUMENT: \{"Váš návrh má celkové úsilie 5419, čo je pod limitom. Pridajme však ešte položku A18, ktorá má úsilie 559, a odstráňme C08 (úsilie 110) a C42 (úsilie 33). Tým získame viac hodnoty."\}} \\
\texttt{NÁVRH: \{'C55', 'A43', 'A67', 'B42', 'A62', 'B19', 'C43', 'C73', 'C10', 'B36', 'A18', 'C42', 'B40', 'C16', 'A21', 'C39'\}} \\
            }
        }
    }
    & \\ \\

    \theutterance \stepcounter{utterance}  
    & & & \multicolumn{4}{p{0.6\linewidth}}{
        \cellcolor[rgb]{0.9,0.9,0.9}{
            \makecell[{{p{\linewidth}}}]{
                \texttt{\tiny{[GM$\langle$P2]}}
                \texttt{ARGUMENT: \{'Položka A18 je už v mojom návrhu zahrnutá. Váš návrh je takmer identický s mojím, len odstraňuje C08. Súhlasím s odstránením C08, ak tým získame väčší súhlas.'\}} \\
\texttt{SÚHLASÍM: \{'C55', 'A43', 'A67', 'B42', 'A62', 'B19', 'C43', 'C73', 'C10', 'B36', 'A18', 'C42', 'B40', 'C16', 'A21', 'C39'\}} \\
            }
        }
    }
     \\ \\

    \theutterance \stepcounter{utterance}  
    & & & \multicolumn{2}{p{0.3\linewidth}}{
        \cellcolor[rgb]{0.9,0.9,0.9}{
            \makecell[{{p{\linewidth}}}]{
                \texttt{\tiny{[GM$|$GM]}}
                \texttt{game successful} \\
            }
        }
    }
    & & \\ \\

    \theutterance \stepcounter{utterance}  
    & & & \multicolumn{2}{p{0.3\linewidth}}{
        \cellcolor[rgb]{0.9,0.9,0.9}{
            \makecell[{{p{\linewidth}}}]{
                \texttt{\tiny{[GM$|$GM]}}
                \texttt{end game} \\
            }
        }
    }
    & & \\ \\

\end{supertabular}
}

\end{document}
